% !TEX program = xelatex
% Отредактированная версия статьи для шаблона НИТУ «МИСИС»
% Основная тема и авторский состав сохранены.

\documentclass{misisarticle}

\usepackage{fontspec}
\setmainfont{Times New Roman}[
  SmallCapsFont = {Times New Roman},
  SmallCapsFeatures = {Letters=SmallCaps, FakeBold=1.2}
]

% --- Классификационные коды ---
\grnti{47.49.31}
\asjc{2208}
\oecd{2.2}

% --- Название ---
\title{Использование UWB-антенн для задач точного позиционирования трамвайных вагонов}

% --- Авторы ---
\misisauthor{М.~А.~Шульгин}{кафедра инженерной кибернетики НИТУ «МИСИС»}{Москва, Россия}{m2311717@edu.misis.ru}

\usepackage{pgfplots}
\usepgfplotslibrary{groupplots}
\pgfplotsset{compat=1.18}
\usepackage{placeins}
\renewcommand{\topfraction}{0.95}
\renewcommand{\textfraction}{0.05}
\renewcommand{\dbltopfraction}{0.95}
\abstract{%
Рассматривается использование UWB-антенн для точного позиционирования трамвайных вагонов на участках с ограниченной доступностью спутниковой навигации. Основное внимание уделено практической проверке дальномерных измерений на двух платах Qorvo DW3000. Одну плату закрепили неподвижно, другую перемещали по прямой с остановками на эталонных отметках. Показания снимали на видео, затем извлекали в общую таблицу. После исключения пустых и нулевых значений проанализированы 925 показаний на 13 расстояниях от 20 до 300 м. Представлены сводные результаты и графики, рассмотрены систематические смещения, разброс и реальные скачки показаний. Показано, что близкая к эталону медиана может сочетаться с большими отклонениями отдельных значений. Полученные результаты характеризуют дальномерный канал испытанной пары плат и позволяют оценить ограничения его использования в задаче позиционирования трамвая. Статистика относится к сохранённым показаниям дисплея и не является оценкой точности независимых радиообменов или координат вагона.}

\keywords{UWB, Qorvo DW3000, дальномерные измерения, систематическое смещение, скачки показаний, видеорегистрация, предварительное исследование.}

\begin{document}
\maketitle

\section{Введение}

Для определения положения трамвая обычно используют спутниковую навигацию. Однако в плотной городской застройке, под перекрытиями и на закрытых участках маршрута условия приёма спутниковых сигналов могут ухудшаться. В таких условиях интерес представляют дополнительные способы определения местоположения, связанные с инфраструктурой пути.

Один из таких способов основан на сверхширокополосной радиосвязи UWB. Радиомодуль на вагоне обменивается сигналами с опорными узлами, координаты которых известны. По измеренным расстояниям и геометрии рельсового пути можно оценивать положение вагона. При этом качество исходных дальномерных измерений напрямую влияет на достоверность такой оценки.

Антенна является частью измерительного тракта: её задержка, ориентация и условия приёма могут влиять на показания дальномера. Для трамвая этот вопрос связан также с металлическим кузовом и изменением направления на опорные узлы при движении. Поэтому в контексте применения UWB-антенн важно понимать не только принцип расчёта координат, но и фактическое поведение радиомодулей, с которых поступают расстояния.

Применение UWB совместно с GPS для локализации транспортного средства в открытых и закрытых пространствах рассмотрено в работе Fernández-Madrigal и соавторов \cite{fernandez2007application}. Queralta и соавторы исследовали UWB-локализацию беспилотников без доступа к ГНСС и представили экспериментальные данные \cite{queralta2020uwb}. Huang и соавторы предложили алгоритм обработки измерений для повышения точности UWB-позиционирования в помещении \cite{huang2024design}. Эти работы показывают, что оценка качества дальномерных данных занимает самостоятельное место в задаче локализации.

Для трамвайного транспорта изучаются и методы на основе изображений и карты: оценка точности на кривых \cite{ali2025} и местоположения в условиях высотной застройки \cite{guzhva2023}. В работе Бикмаева и соавторов рассматривается локализация наземного транспорта по монокамере и дорожным знакам с геодезической привязкой \cite{bikmaev2020}. Исследование UWB относится к той же задаче определения положения транспортного средства, но использует радиотехнические измерения.

В настоящей работе выполнен опыт с двумя платами Qorvo DW3000. Одна плата оставалась неподвижной, вторую перемещали с остановками на известных расстояниях. Цель работы --- оценить, насколько зарегистрированные показания соответствуют контрольным отметкам, и охарактеризовать наблюдавшиеся смещения и скачки. Для этого сформирована выборка остановок, рассчитаны описательные показатели ошибки и построены графики отклонений. Результаты относятся к дальномерным показаниям испытанной пары плат; координаты трамвая в данном опыте не определялись.

\section{Постановка задачи}

\subsection{Измеряемая величина}

Для одностороннего времени распространения $\tau$ после учёта аппаратных задержек расстояние выражается как
\begin{equation}
d=c_0\tau,
\end{equation}
где $c_0$ --- скорость распространения сигнала, приближённо равная скорости света. Документация семейства DW3000 описывает возможности радиотрансивера, включая работу в каналах 5 и 9 \cite{dw3000ds}. Конкретная схема получения дальности зависит от прошивки.

Для остановки на контрольной отметке $d_0$ анализируется отклонение показания $x_k$:
\begin{equation}
e_k=x_k-d_0.
\end{equation}
Значения $d_0$ приняты эталонными для расчёта по разметке опыта. Они не подбираются по дальномеру. При этом способ их измерения, точки отсчёта и инструментальная погрешность не документированы; слово «эталон» здесь обозначает расчётную опору и не утверждает метрологическую прослеживаемость.

\subsection{Границы исследования}

Оцениваются распределение сохранённых показаний, их смещение и крупные отклонения. Последовательность кадров не считается последовательностью независимых радиопакетов. Поэтому не оцениваются доверительные интервалы точности, число независимых обменов, частота обновления дальности или вероятность успешного обмена.

Дальность между двумя узлами и координата вагона --- разные измеряемые величины. В данном опыте зарегистрированы расстояния между платами, а координаты движущегося вагона не определялись. Поэтому их точность по полученной выборке не оценивается.

\section{Методика эксперимента}

\subsection{Стенд и регистрация}

Использованы две платы на базе трансиверов семейства Qorvo DW3000. Одну плату оставляли неподвижной, вторую перемещали по прямой с остановками на контрольных отметках. Выполнялись удаление от опорной платы и обратный проход. Для части измерений на 80 м отмечены препятствия: одна плата находилась на улице, другая --- в помещении.

Показания плат снимали на видео, затем извлекали численные значения в общую таблицу. Отдельно отмечали остановки, расстояния и настройки. Данные рассматриваются как единый набор без описания отдельных видеофайлов. Резкие скачки наблюдались непосредственно на плате; ненулевые скачки сохранены.

Такая регистрация фиксирует отображаемые значения. Она не даёт исходных временных меток DW3000 и не позволяет установить, относится ли очередной кадр к новому измерению. Одинаковые соседние числа могут быть как повторным отображением, так и отдельными результатами. Удалять их автоматически без журнала обменов нельзя.

\subsection{Радионастройки и полнота описания}

Настройки взяты из отдельной колонки журнала (таблица~\ref{tab:settings}). На отметке 300 м изменены канал и ещё два параметра.

\begin{table}[!ht]
\centering
\caption{Параметры в обозначениях интерфейса}
\label{tab:settings}
\begin{tabular}{lcc}
\toprule
Поле & До изменения & После изменения \\
\midrule
Channel & 5 & 9 \\
DR, Мбит/с & 6,8 & 6,8 \\
Gain & 30db & 15db \\
Preamble & 9 & 12 \\
\bottomrule
\end{tabular}
\end{table}

Смысл Gain и Preamble по сохранённой таблице однозначно не установлен. Gain не отождествляется с усилением антенны в dBi или излучаемой мощностью, а Preamble --- с длиной преамбулы. Одновременная смена параметров не позволяет выделить эффект одного канала.

Полный паспорт стенда не сохранён: неизвестны точная модель платы и антенны, версия прошивки, вариант TWR, значения антенных задержек, процедура калибровки, высота и ориентация плат. Поэтому опыт нельзя считать полностью воспроизводимым. В частности, отсутствие сведений о калибровке не доказывает, что её не было; статус индивидуальной калибровки остаётся неподтверждённым. Аппаратные параметры семейства DW3000 описаны в руководстве Qorvo \cite{aph301}; этот справочный документ не заменяет сведений о конкретных экспериментальных платах.

\section{Обработка измерений}

\subsection{Отбор и объединение}

Исходная таблица содержит 2498 строк: в 954 отсутствует извлечённое расстояние, ещё в двух записан явный ноль. Все эти 956 строк исключены по принятому правилу. Остаются 1542 ненулевых показания, включая движение. Для расчёта ошибок отобраны только явно размеченные остановки: 925 значений на 13 расстояниях. Сохранены также 41 значение остановок, перенесённое в итоговое поле из альтернативного результата распознавания.

Пропуски не заполняются, нули не заменяются, сглаживание не применяется. Другие ненулевые значения не отбрасываются по величине ошибки. Уверенность распознавания не интерпретируется как качество радиоканала.

Повторные остановки на одинаковой отметке объединены. Это даёт сводную описательную характеристику, но смешивает условия. На 80 м в выборку входит эпизод с препятствиями, на 300 м --- участок смены настроек. Поэтому строки не являются сериями одного неизменного режима, а сравнение расстояний не устанавливает зависимость ошибки только от дальности.

Нули и пустые строки учитываются в описании отбора, но их доля не объявляется радиопотерями. По видео неизвестны число попыток, причины отсутствия показаний и значение нулевого результата. Все показатели ниже условны: они относятся только к сохранённым ненулевым показаниям остановок. Доступность системы по этим данным не оценена.

\subsection{Описательные показатели}

Для $N$ сохранённых значений на отметке $d_0$ рассчитываются
\begin{align}
\bar{x}&=\frac{1}{N}\sum_{k=1}^{N}x_k,\\
b&=\bar{x}-d_0,\\
\mathrm{MAE}&=\frac{1}{N}\sum_{k=1}^{N}|x_k-d_0|,\\
s_x&=\sqrt{\frac{1}{N-1}\sum_{k=1}^{N}(x_k-\bar{x})^2}.
\end{align}
Здесь $\bar{x}$ --- среднее показание, $b$ --- среднее отклонение от отметки, $s_x$ --- выборочный разброс сохранённых значений. Также приводятся медиана и 95-й процентиль абсолютного отклонения $P_{95}$. Для него применяется линейная интерполяция упорядоченной выборки с индексом $(N-1)\cdot0{,}95$ при нумерации с нуля.

$N$ обозначает число извлечённых значений, а не число независимых обменов. При повторении дисплея длительно отображавшийся результат получает больший вес. Поэтому $s_x$ и $P_{95}$ описывают сохранённую выборку, но не задают распределение ошибок независимых измерений. Округление до сотых метра не означает соответствующую достоверность эталона или аппаратуры.

\section{Результаты}

Сводная статистика приведена в таблице~\ref{tab:results}. Внутри каждой строки объединены все остановки на соответствующей отметке. На рисунке~\ref{fig:errors} показаны отклонения в порядке сохранения данных. Ось номера не является временной шкалой: участки движения и исключённые значения в ней отсутствуют.

\begin{table*}[!t]
\centering
\caption{Описательная статистика 925 сохранённых ненулевых показаний остановок. Все величины, кроме $N$, в метрах. Повторы дисплея возможны; условия объединены, включая препятствия на 80 м и смену конфигурации на 300 м}
\label{tab:results}
\begin{tabular}{rrrrrrrr}
\toprule
$d_0$ & $N$ & $\bar{x}$ & Медиана & $b$ & MAE & $s_x$ & $P_{95}$ \\
\midrule
20 & 50 & 20,86 & 20,95 & 0,86 & 0,86 & 0,38 & 1,31 \\
25 & 30 & 26,06 & 26,10 & 1,06 & 1,06 & 0,46 & 1,48 \\
60 & 72 & 61,65 & 61,64 & 1,65 & 1,67 & 0,44 & 2,09 \\
80 & 56 & 86,31 & 84,56 & 6,31 & 6,32 & 4,55 & 14,84 \\
95 & 91 & 96,63 & 95,80 & 1,63 & 1,80 & 2,78 & 7,16 \\
130 & 162 & 139,23 & 131,03 & 9,23 & 17,92 & 40,83 & 75,94 \\
180 & 84 & 258,62 & 180,98 & 78,62 & 79,66 & 93,87 & 215,54 \\
185 & 107 & 193,98 & 184,52 & 8,98 & 17,20 & 36,47 & 107,21 \\
210 & 38 & 242,84 & 213,89 & 32,84 & 68,91 & 78,03 & 113,53 \\
223 & 97 & 243,68 & 223,83 & 20,68 & 21,07 & 41,67 & 128,91 \\
256 & 36 & 256,66 & 256,75 & 0,66 & 0,66 & 0,25 & 0,89 \\
259 & 77 & 288,88 & 260,08 & 29,88 & 30,09 & 56,62 & 141,31 \\
300 & 25 & 300,44 & 300,40 & 0,44 & 0,44 & 0,28 & 0,94 \\
\bottomrule
\end{tabular}
\end{table*}

\subsection{Смещение на ближних отметках}

На 20, 25 и 60 м среднее отклонение положительно и равно соответственно 0,86, 1,06 и 1,65 м. Выборочное стандартное отклонение составляет 0,38--0,46 м. Таким образом, относительно компактная выборка сохраняет заметное завышение дальности. Таким образом, положительное смещение является наблюдаемым свойством этой выборки; его источник по имеющимся данным не установлен.

\subsection{Крупные отклонения}

На 130 м медиана равна 131,03 м, а MAE --- 17,92 м; значения лежат от 62,66 до 351,27 м. На 180 м медиана 180,98 м соседствует с MAE 79,66 м и максимальным показанием 395,89 м. На этой отметке 39 из 84 сохранённых значений имеют абсолютное отклонение более 10 м. Это доля внутри отобранной выборки, а не вероятность ошибки независимого обмена.

Различие медианы и MAE показано на рисунке~\ref{fig:errors}. Описательно оно указывает на значительный вклад крайних значений. Параметрическая модель распределения и наличие «тяжёлого хвоста» в строгом статистическом смысле здесь не проверяются.

\subsection{Участки с малым разбросом и смешанные условия}

На 256 м среднее показание составляет 256,66 м, MAE --- 0,66 м, $s_x$ --- 0,25 м. На 300 м сохранены 25 значений со средним 300,44 м и MAE 0,44 м. Вторая выборка относится к участку изменения конфигурации. Эти числа характеризуют конкретные записи, но не доказывают, что большая дальность улучшает результат или что канал 9 точнее канала 5.

На 80 м общая выборка содержит 56 значений, включая эпизод с препятствиями. Среднее равно 86,31 м, MAE --- 6,32 м. Отдельный эффект препятствия не оценивается: условия до и после его появления не сведены к сопоставимому контрольному опыту.

\begin{figure*}[!t]
\centering
\begin{tikzpicture}
\begin{groupplot}[group style={group size=1 by 2,vertical sep=1.35cm},
width=0.94\textwidth,height=5.2cm,
tick label style={font=\small},label style={font=\small},
grid=major,major grid style={gray!20},scaled ticks=false]
\nextgroupplot[xmin=0,xmax=930,ymin=-100,ymax=260,
xlabel={Номер сохранённого показания},ylabel={Отклонение, м}]
\addplot[only marks,mark=*,mark size=0.75pt,blue!70!black] coordinates {(1,0.0100) (2,0.2300) (3,0.2700) (4,0.3300) (5,0.3200) (6,0.3200) (7,0.3300) (8,0.3400) (9,0.3500) (10,0.3500) (11,0.3600) (12,0.6000) (13,0.6100) (14,1.3100) (15,1.3100) (16,1.1400) (17,1.3000) (18,1.0900) (19,0.9800) (20,1.0000) (21,1.0200) (22,1.1500) (23,1.1600) (24,1.2500) (25,1.2500) (26,1.2700) (27,1.2800) (28,1.2800) (29,1.2100) (30,1.0700) (31,1.0900) (32,0.9600) (33,0.8500) (34,0.8300) (35,0.7000) (36,1.5500) (37,1.3300) (38,1.2100) (39,1.1100) (40,0.9400) (41,0.9400) (42,0.8300) (43,0.7000) (44,0.6000) (45,0.4900) (46,1.2000) (47,1.0300) (48,0.7800) (49,0.6700) (50,0.5700) (51,-0.5600) (52,1.2200) (53,1.0200) (54,1.6900) (55,1.7200) (56,2.1700) (57,1.8100) (58,2.0800) (59,1.7500) (60,1.8900) (61,1.7600) (62,1.6300) (63,1.8800) (64,1.8800) (65,1.7100) (66,1.7600) (67,1.8300) (68,2.8200) (69,2.0100) (70,2.0600) (71,2.0200) (72,2.0200) (73,2.1000) (74,1.3300) (75,1.4500) (76,1.4600) (77,1.4700) (78,1.4700) (79,1.4700) (80,1.4700) (81,1.4800) (82,1.5900) (83,1.4600) (84,3.1600) (85,1.8800) (86,1.8600) (87,1.8700) (88,1.8900) (89,1.8900) (90,1.8000) (91,1.7800) (92,1.7900) (93,1.6800) (94,1.8700) (95,1.6400) (96,1.6200) (97,1.6100) (98,1.6000) (99,1.5900) (100,1.5900) (101,1.5800) (102,1.6100) (103,1.6200) (104,1.4900) (105,1.5000) (106,1.7000) (107,1.5200) (108,1.1300) (109,1.1100) (110,1.2900) (111,1.3100) (112,1.3000) (113,1.3100) (114,1.3100) (115,1.0800) (116,1.3400) (117,1.9000) (118,1.2200) (119,1.6300) (120,2.0800) (121,2.0700) (122,1.7100) (123,-3.1700) (124,-1.8900) (125,-1.4700) (126,-0.5200) (127,-0.3100) (128,-0.0100) (129,0.0100) (130,0.1100) (131,0.0800) (132,-0.1500) (133,-0.0200) (134,-0.0200) (135,0.5700) (136,0.4700) (137,0.5900) (138,-0.0900) (139,0.0700) (140,0.2700) (141,0.2000) (142,0.7800) (143,0.6000) (144,0.6000) (145,0.7700) (146,-0.1100) (147,-0.0400) (148,0.0100) (149,-0.0100) (150,0.7300) (151,0.4400) (152,0.3000) (153,0.4400) (154,0.4500) (155,0.4300) (156,0.4200) (157,1.0900) (158,1.0900) (159,1.0900) (160,1.0900) (161,1.1300) (162,0.9900) (163,0.8500) (164,0.9000) (165,0.8900) (166,0.7300) (167,0.7300) (168,0.8400) (169,0.8500) (170,0.8500) (171,0.8100) (172,0.6300) (173,0.3000) (174,0.3000) (175,0.2900) (176,0.2900) (177,0.2900) (178,0.2800) (179,0.8100) (180,0.8000) (181,1.9400) (182,1.6100) (183,1.6200) (184,0.7600) (185,0.6900) (186,0.9800) (187,0.9800) (188,0.9700) (189,0.9700) (190,2.9900) (191,3.3400) (192,3.0600) (193,5.3400) (194,4.6900) (195,7.0300) (196,7.2800) (197,6.8200) (198,5.0400) (199,4.3200) (200,3.9100) (201,2.7600) (202,2.4000) (203,1.9800) (204,13.5800) (205,13.5700) (206,0.9700) (207,0.9600) (208,7.7400) (209,1.7600) (210,3.8500) (211,3.8600) (212,10.6800) (213,1.1500) (214,-8.6600) (215,-3.7600) (216,-2.5000) (217,-27.0500) (218,-38.3890) (219,-0.3200) (220,0.3700) (221,5.4200) (222,-14.5800) (223,-27.3500) (224,-27.3100) (225,-26.1700) (226,-27.2600) (227,-45.4200) (228,5.5900) (229,-28.5100) (230,3.1600) (231,2.0900) (232,1.4700) (233,-19.2000) (234,-1.2300) (235,-1.7800) (236,-2.2600) (237,-2.4000) (238,-2.4000) (239,-2.2800) (240,-22.0600) (241,-2.0800) (242,2.0600) (243,-2.0800) (244,-2.0900) (245,-2.1100) (246,5.6900) (247,13.5000) (248,13.2100) (249,13.0300) (250,13.0300) (251,13.0200) (252,12.9400) (253,-27.2800) (254,-67.3400) (255,12.9100) (256,12.9300) (257,-4.6000) (258,12.9400) (259,12.9400) (260,13.0100) (261,12.9200) (262,13.1500) (263,13.2100) (264,13.2700) (265,13.4100) (266,13.4000) (267,13.3800) (268,13.2900) (269,13.4000) (270,13.2500) (271,13.1000) (272,12.9700) (273,12.7300) (274,12.7400) (275,12.7600) (276,13.0100) (277,13.0500) (278,13.0300) (279,13.6900) (280,12.9400) (281,-2.4800) (282,-2.6100) (283,-2.7500) (284,-0.1200) (285,0.7000) (286,0.7000) (287,3.5700) (288,-2.2200) (289,-2.3000) (290,-2.5700) (291,0.6700) (292,-2.8100) (293,-49.8400) (294,-2.1200) (295,-2.3300) (296,-2.9700) (297,-2.1500) (298,-3.0300) (299,1.0700) (300,10.5000) (301,6.2600) (302,-4.0900) (303,-4.5500) (304,-5.1400) (305,-5.6200) (306,-7.2600) (307,-8.1200) (308,-8.7400) (309,110.0000) (310,-10.3300) (311,-11.0000) (312,5.5700) (313,-12.2400) (314,-13.7200) (315,-14.9900) (316,-3.1600) (317,-16.8700) (318,179.6000) (319,221.2700) (320,178.8800) (321,30.1700) (322,16.5700) (323,217.5000) (324,216.9700) (325,76.3900) (326,-19.7000) (327,-19.0200) (328,-17.8300) (329,-17.1800) (330,-16.4200) (331,-15.5500) (332,107.2400) (333,-14.1000) (334,86.5900) (335,70.1700) (336,70.1300) (337,124.2000) (338,70.1300) (339,-10.6300) (340,-8.8100) (341,-8.2100) (342,-7.3900) (343,-5.7500) (344,70.8400) (345,111.6300) (346,-8.1200) (347,143.0600) (348,108.6300) (349,141.3300) (350,-9.2300) (351,107.1300) (352,-8.6600) (353,-8.5500) (354,-8.6500) (355,-8.6400) (356,-8.6400) (357,-8.6300) (358,-8.6100) (359,-8.6100) (360,-8.3700) (361,-0.5100) (362,-0.5000) (363,-0.5100) (364,-0.6300) (365,-0.6100) (366,-0.6300) (367,-0.5200) (368,-0.5000) (369,-0.4800) (370,-0.4500) (371,-0.0400) (372,-8.7800) (373,-8.7700) (374,-8.5700) (375,-8.6800) (376,-8.6800) (377,-8.7900) (378,-9.0800) (379,-11.3500) (380,-10.1300) (381,-10.7500) (382,-10.7500) (383,-11.7000) (384,0.1400) (385,0.1200) (386,-0.0200) (387,-0.0100) (388,0.0000) (389,0.0200) (390,0.1200) (391,0.1000) (392,0.1000) (393,0.1100) (394,0.1000) (395,-0.0100) (396,0.0100) (397,-0.4800) (398,-0.4700) (399,-0.5800) (400,-0.6000) (401,-0.2200) (402,-0.1800) (403,0.1900) (404,0.1400) (405,-4.2700) (406,-8.6800) (407,-8.7100) (408,-6.7200) (409,1.3200) (410,0.0400) (411,0.0200) (412,0.0400) (413,0.0900) (414,-2.3900) (415,-0.1700) (416,-0.3100) (417,-0.3100) (418,0.1100) (419,-10.1500) (420,21.9900) (421,65.2900) (422,97.8400) (423,-3.7000) (424,0.3400) (425,0.3300) (426,0.3500) (427,0.2000) (428,0.2000) (429,0.1800) (430,0.1700) (431,0.1700) (432,-15.1900) (433,111.3700) (434,112.1200) (435,111.4600) (436,110.8000) (437,110.3100) (438,109.7800) (439,109.2700) (440,108.5300) (441,107.9800) (442,80.1400) (443,106.7100) (444,106.2200) (445,-40.8600) (446,-43.5500) (447,-48.2000) (448,-52.3500) (449,-53.0500) (450,-50.5000) (451,-51.0600) (452,-51.8500) (453,-7.7700) (454,-6.4400) (455,-5.7800) (456,-5.2700) (457,-60.6400) (458,40.2000) (459,80.1100) (460,89.2200) (461,-0.2800) (462,-59.8200) (463,0.9200) (464,-62.0000) (465,-32.2400) (466,69.3700) (467,121.4900) (468,240.3900) (469,-53.7400) (470,6.8500) (471,0.1600) (472,0.3700) (473,0.3600) (474,0.3800) (475,0.3800) (476,0.3600) (477,0.2100) (478,0.1900) (479,0.2000) (480,0.2000) (481,0.2100) (482,0.2200) (483,0.2000) (484,0.1900) (485,0.2000) (486,0.2000) (487,0.4900) (488,0.4900) (489,0.4900) (490,0.4800) (491,0.4800) (492,0.6200) (493,141.3100) (494,141.3200) (495,141.3200) (496,141.3200) (497,141.3200) (498,141.3100) (499,1.4100) (500,1.4100) (501,1.3000) (502,1.2400) (503,1.2200) (504,1.2200) (505,141.2700) (506,31.8600) (507,110.6000) (508,1.4200) (509,141.2400) (510,141.2400) (511,141.2400) (512,141.2400) (513,141.2400) (514,1.2700) (515,1.4000) (516,1.4000) (517,1.3900) (518,1.2500) (519,1.1200) (520,-8.0800) (521,1.0300) (522,0.9300) (523,1.0800) (524,1.0800) (525,1.0700) (526,1.0600) (527,0.9300) (528,0.8100) (529,0.6700) (530,0.6600) (531,0.6400) (532,0.6400) (533,0.6300) (534,0.6200) (535,0.6200) (536,1.0400) (537,1.1500) (538,1.4400) (539,1.5800) (540,140.8300) (541,140.8000) (542,1.4200) (543,140.8200) (544,0.7100) (545,1.3900) (546,1.3900) (547,1.6500) (548,0.5700) (549,0.5700) (550,0.5500) (551,0.4100) (552,0.3200) (553,0.1500) (554,0.1400) (555,0.1300) (556,0.1200) (557,0.2500) (558,0.4100) (559,0.4000) (560,0.4000) (561,0.5100) (562,0.4800) (563,0.7600) (564,0.8600) (565,0.9800) (566,0.9600) (567,0.8600) (568,0.3300) (569,0.1700) (570,0.0400) (571,0.3200) (572,0.1900) (573,0.9900) (574,0.6000) (575,0.4800) (576,0.3900) (577,0.6300) (578,0.8100) (579,0.8200) (580,0.8100) (581,0.8000) (582,0.7900) (583,0.7900) (584,0.7900) (585,0.3700) (586,0.3700) (587,0.3400) (588,0.6900) (589,0.6800) (590,0.6600) (591,0.6400) (592,0.6300) (593,0.6200) (594,0.7200) (595,0.8700) (596,0.8600) (597,0.8500) (598,0.8400) (599,0.8400) (600,0.8400) (601,0.8500) (602,0.8500) (603,0.9400) (604,0.7800) (605,0.5500) (606,0.1600) (607,-0.0100) (608,-0.0300) (609,0.7600) (610,0.5600) (611,0.7200) (612,0.8700) (613,0.8700) (614,0.8500) (615,0.8300) (616,0.8300) (617,0.8300) (618,0.8300) (619,0.8400) (620,0.8400) (621,0.9900) (622,1.1000) (623,65.2000) (624,1.2500) (625,1.2200) (626,65.0000) (627,32.4000) (628,1.0400) (629,1.0800) (630,1.0800) (631,1.0500) (632,0.7800) (633,0.7900) (634,0.7800) (635,0.7700) (636,0.5900) (637,0.4400) (638,0.2200) (639,0.0900) (640,0.0900) (641,0.1000) (642,0.1100) (643,0.1100) (644,0.1100) (645,65.1100) (646,28.8300) (647,0.9000) (648,0.5300) (649,0.5200) (650,0.5100) (651,0.5100) (652,0.5000) (653,0.3600) (654,0.7400) (655,0.7300) (656,0.7900) (657,0.7800) (658,0.7700) (659,0.7600) (660,0.7500) (661,0.7600) (662,0.7600) (663,29.1000) (664,0.8600) (665,0.8600) (666,0.8500) (667,0.8300) (668,0.8100) (669,1.8500) (670,96.4400) (671,32.4900) (672,96.4200) (673,129.0000) (674,114.7000) (675,128.9800) (676,128.9900) (677,101.1300) (678,1.3800) (679,0.8900) (680,0.7600) (681,0.7400) (682,0.7300) (683,0.7100) (684,0.7100) (685,128.9000) (686,128.9600) (687,128.9100) (688,128.9100) (689,128.8900) (690,37.4100) (691,1.1200) (692,61.3500) (693,1.4200) (694,6.4400) (695,1.2800) (696,105.1800) (697,1.3900) (698,-19.1200) (699,0.8000) (700,0.7500) (701,0.7600) (702,0.7700) (703,1.1100) (704,0.3900) (705,0.0400) (706,0.1500) (707,-0.1600) (708,-0.0100) (709,-0.0200) (710,75.5200) (711,75.5400) (712,107.7800) (713,33.4400) (714,150.6400) (715,150.7500) (716,207.0600) (717,150.7200) (718,0.6200) (719,117.8700) (720,0.6200) (721,108.4800) (722,-3.4300) (723,195.5400) (724,0.1000) (725,0.0200) (726,0.2500) (727,0.2000) (728,0.0000) (729,-0.3300) (730,-0.3100) (731,-0.3100) (732,-0.0600) (733,0.1000) (734,0.2400) (735,0.2600) (736,0.2400) (737,0.2300) (738,0.6500) (739,168.4800) (740,-4.6600) (741,57.2800) (742,55.3900) (743,215.4800) (744,167.3100) (745,215.4600) (746,215.4700) (747,215.4700) (748,-4.6700) (749,0.2300) (750,0.2300) (751,0.2300) (752,0.9400) (753,8.5500) (754,0.3300) (755,-0.1400) (756,-0.3800) (757,-2.5300) (758,0.9800) (759,0.9800) (760,0.9800) (761,215.5500) (762,215.8900) (763,215.5000) (764,215.5200) (765,215.5400) (766,215.5700) (767,198.2400) (768,215.5700) (769,213.5800) (770,215.5200) (771,105.4200) (772,215.5100) (773,215.5000) (774,215.5000) (775,105.3200) (776,215.4300) (777,-4.8200) (778,-4.6800) (779,215.4300) (780,5.1800) (781,215.5400) (782,150.3100) (783,150.2800) (784,-4.7900) (785,-3.1100) (786,-4.7300) (787,-4.7100) (788,0.6700) (789,0.3200) (790,2.5900) (791,10.2500) (792,0.9900) (793,1.0300) (794,0.9500) (795,1.0300) (796,0.8900) (797,1.0700) (798,1.0800) (799,1.1300) (800,9.9000) (801,9.8800) (802,9.8700) (803,9.8600) (804,9.8600) (805,9.8600) (806,0.6800) (807,0.7000) (808,0.7800) (809,0.6900) (810,0.5000) (811,0.4600) (812,-0.6200) (813,116.7200) (814,116.7200) (815,-3.1400) (816,-2.8400) (817,-2.8600) (818,-2.7700) (819,-2.8000) (820,-2.9400) (821,13.4600) (822,13.4500) (823,13.4500) (824,13.4400) (825,13.4400) (826,13.4400) (827,13.4500) (828,13.4500) (829,13.4500) (830,13.6000) (831,5.7800) (832,-35.5100) (833,-2.8900) (834,-2.7400) (835,13.3300) (836,8.6500) (837,7.2800) (838,3.2600) (839,0.6300) (840,1.7600) (841,1.6200) (842,1.7300) (843,1.6800) (844,1.9500) (845,2.1100) (846,2.2800) (847,2.1100) (848,2.2200) (849,2.1000) (850,1.3700) (851,2.6200) (852,2.2200) (853,2.1000) (854,2.7200) (855,2.5900) (856,1.9000) (857,2.4500) (858,2.5600) (859,2.6600) (860,4.2600) (861,2.7500) (862,2.8200) (863,6.3900) (864,8.0600) (865,9.7600) (866,8.5400) (867,8.9900) (868,8.9900) (869,10.3500) (870,10.5800) (871,10.7600) (872,10.6300) (873,10.1400) (874,10.1500) (875,10.1500) (876,10.2400) (877,14.8500) (878,15.0300) (879,13.6900) (880,14.8400) (881,14.8400) (882,13.8600) (883,13.4600) (884,10.7300) (885,10.0200) (886,9.4500) (887,8.0300) (888,5.8700) (889,5.3000) (890,4.8500) (891,4.1800) (892,3.5200) (893,2.9700) (894,2.0100) (895,-0.2400) (896,1.0400) (897,1.2600) (898,0.7300) (899,0.2200) (900,0.4700) (901,0.6200) (902,0.7100) (903,1.0900) (904,1.1300) (905,1.1100) (906,0.8800) (907,0.8700) (908,0.9000) (909,1.1000) (910,2.7800) (911,1.1300) (912,0.7900) (913,1.4700) (914,1.3100) (915,1.2900) (916,1.4800) (917,1.1800) (918,0.9500) (919,1.4300) (920,1.1000) (921,1.1200) (922,1.1400) (923,1.2100) (924,0.9800) (925,0.1600)};
\addplot[gray,dashed,domain=0:930] {0};
\nextgroupplot[ymode=log,ymin=0.1,ymax=300,
xmin=0.6,xmax=13.4,xtick={1,2,3,4,5,6,7,8,9,10,11,12,13},
xticklabels={20,25,60,80,95,130,180,185,210,223,256,259,300},
xlabel={Контрольное расстояние, м},ylabel={Отклонение, м (лог.)},
legend style={at={(0.02,0.98)},anchor=north west,legend columns=2,font=\small}]
\addplot[orange!80!black,thick,mark=square*] coordinates {(1,0.85740) (2,1.05500) (3,1.66625) (4,6.32232) (5,1.79857) (6,17.92073) (7,79.66131) (8,17.20364) (9,68.91158) (10,21.06969) (11,0.65806) (12,30.08740) (13,0.43520)};
\addlegendentry{MAE}
\addplot[blue!70!black,thick,mark=*] coordinates {(1,0.95000) (2,1.10000) (3,1.63500) (4,4.55500) (5,0.80000) (6,1.03000) (7,0.98000) (8,0.48000) (9,3.88500) (10,0.83000) (11,0.75000) (12,1.08000) (13,0.40000)};
\addlegendentry{$|\mathrm{median}(x)-d_0|$}
\end{groupplot}
\end{tikzpicture}
\caption{Отклонения сохранённых ненулевых показаний остановок: сверху --- в порядке регистрации, снизу --- MAE и абсолютное отклонение медианы по контрольным расстояниям. Пропуски и движение не показаны, ось номера не задаёт время. Нижняя шкала логарифмическая. Все условия объединены, включая препятствия на 80 м и изменение настроек на 300 м.}
\label{fig:errors}
\end{figure*}

\FloatBarrier

\section{Обсуждение}

Основной результат опыта --- сочетание положительного смещения на ближних отметках и крупных скачков на части более дальних участков. Малый разброс сам по себе не означает близость к контрольному расстоянию, а близкая к нему медиана не исключает больших отклонений. Этот вывод следует из сохранённых показаний без сглаживания и удаления ненулевых скачков.

Причину скачков по таблице расстояний установить нельзя. В частности, влияние среды, временных оценок и прошивки не разделено. Наблюдения не позволяют приписать их антенне или металлическому кузову: антенна как отдельный фактор и вагон в опыте не исследовались. Не определены также оптимальное место монтажа и шаг установки опорных узлов.

Неизвестная погрешность контрольных отметок ограничивает метрологическую трактовку смещения. Если фактическое расстояние отличается от принятого $d_0$ на $\Delta d_0$, рассчитанное среднее отклонение изменится на $-\Delta d_0$. Без оценки $\Delta d_0$ нельзя считать MAE 0,44 м на 300 м подтверждённой абсолютной точностью дальномера. Однако сама нестабильность показаний на размеченной остановке не устраняется небольшой постоянной поправкой контрольной дистанции.

Повторы дисплея, исключение нулей и пропусков и объединение разных режимов дополнительно ограничивают обобщение. Отсутствие сведений обо всех попытках обмена и их статусах не позволяет оценить доступность системы. Сохранённое показание нельзя приравнять к успешному независимому обмену, а отсутствие показания --- к потере пакета.

В контексте позиционирования трамвая результаты показывают ограничения исходных дальномерных данных: положительное смещение и крупные скачки влияют на достоверность расстояний, используемых для расчёта положения. Гарантированная дальность, точность координат и пригодность для задач безопасности по этому опыту не устанавливаются.

\section{Заключение}

В предварительном опыте с парой плат DW3000 сформирована выборка из 925 сохранённых ненулевых показаний остановок на 13 контрольных расстояниях. На 20--60 м зарегистрировано завышение среднего на 0,86--1,65 м при выборочном разбросе менее 0,5 м. На части более дальних участков выявлены крупные скачки; на 180 м MAE составила 79,66 м при медиане 180,98 м.

Получены описательная статистика и графическое представление отклонений. Они показывают, что оценки только по среднему или медиане недостаточно для характеристики зарегистрированных показаний. При этом точность независимых радиообменов, доля успешных измерений, влияние антенного монтажа и точность координат трамвая не определены.

\begin{thebibliography}{99}

\bibitem{fernandez2007application}
J.~A.~Fernández-Madrigal, E.~Cruz-Martín, J.~González, C.~Galindo, and J.~L.~Blanco,
“Application of UWB and GPS technologies for vehicle localization in combined indoor-outdoor environments,”
in \textit{2007 9th International Symposium on Signal Processing and Its Applications},
Sharjah, United Arab Emirates, 2007, pp.~1--4.
DOI: \href{https://doi.org/10.1109/ISSPA.2007.4555416}{10.1109/ISSPA.2007.4555416}.

\bibitem{huang2024design}
Y.~Huang, B.~Cao, and A.~Wang,
“Design a novel algorithm for enhancing UWB positioning accuracy in GPS denied environments,”
\textit{Scientific Reports}, vol.~14, Art.~no.~23895, 2024.
DOI: \href{https://doi.org/10.1038/s41598-024-74773-y}{10.1038/s41598-024-74773-y}.

\bibitem{queralta2020uwb}
J.~P.~Queralta, C.~Martínez Almansa, F.~Schiano, D.~Floreano, and T.~Westerlund,
“UWB-based System for UAV Localization in GNSS-Denied Environments: Characterization and Dataset,”
in \textit{2020 IEEE/RSJ International Conference on Intelligent Robots and Systems (IROS)},
Las Vegas, NV, USA, 2020, pp.~4521--4528.
DOI: \href{https://doi.org/10.1109/IROS45743.2020.9341042}{10.1109/IROS45743.2020.9341042}.

\bibitem{ali2025}
Б.~Али, Р.~Н.~Садеков,
«Оценка точности позиционирования трамвая в кривых на основе данных карты и сегментированных изображений»,
\textit{Научно-технический вестник информационных технологий, механики и оптики},
2025, т.~25, №~4, с.~771--779.
DOI: \href{https://doi.org/10.17586/2226-1494-2025-25-4-771-779}{10.17586/2226-1494-2025-25-4-771-779}.

\bibitem{guzhva2023}
Н.~С.~Гужва, Б.~Али, К.~С.~Бакулев и др.,
«Оценка точности местоположения трамвая в условиях высотной застройки по данным визуально-геоинформационных систем»,
в: \textit{XXX Юбилейная Санкт-Петербургская международная конференция по интегрированным навигационным системам: сборник материалов конференции},
Санкт-Петербург, 29--31 мая 2023~г.
Санкт-Петербург: Концерн «ЦНИИ „Электроприбор“», 2023, с.~16--21.
EDN: VCHMDX.

\bibitem{bikmaev2020}
Р.~Р.~Бикмаев, А.~А.~Полукаров, Р.~Н.~Садеков,
«Определение местоположения наземного транспортного средства с использованием монокамеры и дорожных знаков с геодезической привязкой»,
в: \textit{XXVII Санкт-Петербургская международная конференция по интегрированным навигационным системам: сборник материалов},
Санкт-Петербург, 25 мая -- 5 июня 2020~г.
Санкт-Петербург: Концерн «ЦНИИ „Электроприбор“», 2020, с.~341--345.
EDN: TDXGEK.



\bibitem{dw3000ds}
Qorvo. \textit{DW3000 family: Product Data Sheet}, Rev.~1.3, 2020.
Документ в разделе Technical Documents страницы DW3110.
\href{https://www.qorvo.com/products/p/DW3110}{qorvo.com/products/p/DW3110}.
Семейство используется как справочный источник, точная модель экспериментальной платы не установлена.

\bibitem{aph301}
Qorvo. \textit{Hardware Design Guide for DW3000 and QM33100 Series ICs (APH301)}, Rev.~A, January 2024.
Раздел Technical Documents:
\href{https://www.qorvo.com/products/p/DW3110}{qorvo.com/products/p/DW3110}.

\end{thebibliography}
\end{document}
